\exercisesection{扩张模}

¶ 1) 一个 A-复形 C 被称为纯的，如果对每个右 A-模 P，有 H(P ⊗\_A C) = 0。

\begin{enumerate}
\def\labelenumi{\alph{enumi})}
\item
  同伦于零的复形是纯的；由平坦模构成、同调为零且有右界的复形是纯的。给出一个秩为 1 的自由模复形的例子，其同调为零但不纯。
\item
  证明以下条件是等价的：
\end{enumerate}

α) C 是纯的。

β) 对每个有限表示的 A-模 M，有 H(Homgr\_A (M, C)) = 0。

\(\gamma)\) 存在一个滤过的归纳系统 \((C^{(\alpha)})_{\alpha \in j}\) 由同伦于零的复形组成，使得 \(C = \varinjlim C^{(\alpha)}\) .

(为了看出 \(\beta\) ）蕴含 \(\gamma\) )，化简到 \(C_i = 0\) 当 \(i < i_0\) 的情形，并将 \(C_{i_0}\) 写成有限表示模的归纳极限。

\begin{enumerate}
\def\labelenumi{\arabic{enumi})}
\setcounter{enumi}{1}
\tightlist
\item
  设 M 是一个 A-模。我们称 M 的子模 \(\mathbf{M}'\) 是纯的，如果由正合序列 \(0 \to \mathbf{M}' \to \mathbf{M} \to \mathbf{M} / \mathbf{M}' \to 0\) 所定义的复形是纯的（练习 1）。
\end{enumerate}

\begin{enumerate}
\def\labelenumi{\alph{enumi})}
\item
  证明：当A是主理想环时，纯子模的概念与VII，§2，习题7中的概念一致。
\item
  假设A-模M是平坦的。证明以下条件是等价的： \(\alpha\) 子模 M' 是纯的；
\end{enumerate}

β) 商模 M/M' 是平坦的；

\(\gamma)\) 对每个理想 \(\mathfrak{a}\) ⊂ \(\mathbf{A}\) ，都有 \(\mathbf{M}' \cap \mathfrak{a}\mathbf{M} = \mathfrak{a}\mathbf{M}'\) 。

\begin{enumerate}
\def\labelenumi{\arabic{enumi})}
\setcounter{enumi}{2}
\tightlist
\item
  设 M 是一个 A-模，n 为整数且 n ≥ 1。证明以下条件等价：
\end{enumerate}

α) M 具有有限 n-表示（X，第 180 页，习题 6）。

β) 对于每个滤过归纳系统 \((\mathrm{N}_{\alpha})_{\alpha\in I}\) ，典范同态

\[
\varinjlim \operatorname{Ext} _ {\mathbf {A}} ^ {i} (\mathbf {M}, \mathbf {N} _ {\alpha}) \rightarrow \operatorname{Ext} _ {\mathbf {A}} ^ {i} (\mathbf {M}, \varinjlim \mathbf {N} _ {\alpha})
\]

对 i \textless{} n 是双射的。

γ) 对于右 A-模的每一个族 \((\mathbf{P}_{\beta})_{\beta \in J}\) ，典范同态

\[
\operatorname{Tor} _ {i} ^ {\mathbf {A}} \left(\prod_ {\beta \in J} \mathrm{P} _ {\beta}, \mathrm{M}\right)\rightarrow \prod_ {\beta \in J} \operatorname{Tor} _ {i} ^ {\mathbf {A}} \left(\mathrm{P} _ {\beta}, \mathrm{M}\right) \quad \text { is   bijective   for } \quad i <   n.
\]

δ) 对于每个集合 I，典范同态 \(A_{d}^{I} \otimes_{A} M \to M^{I}\) 是双射的，并且我们有

\[
\operatorname{Tor} _ {i} ^ {\mathbf {A}} \left(\mathrm{A} _ {d} ^ {\mathrm{I}}, \mathrm{M}\right) = 0 \quad \text { for } \quad 0 <   i <   n.
\]

（对 n 作归纳论证，利用第 170 页练习 18。）

\begin{enumerate}
\def\labelenumi{\arabic{enumi})}
\setcounter{enumi}{3}
\tightlist
\item
  设 P 为右零的投射 A-模复形，且 n 为满足 \(H_{i}(P)=0\) （i\textless n ）的整数。设 M 为左 A-模，N 为右 A-模。证明我们有 \(H^{i}(\mathrm{Homgr}_{\mathbf{A}}(\mathbf{P},\mathbf{M}))=0\) 和 \(H_{i}(\mathbf{N}\otimes_{\mathbf{A}}\mathbf{P})=0\) （i\textless n ），且典范同态
\end{enumerate}

\[
\lambda^ {n}: \mathrm{H} ^ {n} (\operatorname{Homgr} _ {\mathbf {A}} (\mathrm{P}, \mathrm{M})) \rightarrow \operatorname{Hom} _ {\mathbf {A}} \left(\mathrm{H} _ {n} (\mathrm{P}), \mathrm{M}\right), \quad \gamma^ {n}: \mathrm{N} \otimes_ {\mathbf {A}} \mathrm{H} _ {n} (\mathrm{P}) \rightarrow \mathrm{H} _ {n} (\mathrm{N} \otimes_ {\mathbf {A}} \mathrm{P})
\]

都是双射的。

¶ 5) 设 P 是一个右零的投射 A-复形，n 是一个整数 ≥ 0。证明以下条件是等价的：

α) 存在一个由有限生成投射 A-模构成的复形 \(\mathbf{P}'\) ，使得 \(\mathbf{P}_i' = 0\) （对 \(i > n\) 或 \(i < 0\) ），以及一个态射 \(f: \mathbf{P}' \to \mathbf{P}\) ，使得 \(\mathbf{H}_i(f)\) 对 \(i < n\) 为双射且对 \(i = n\) 为满射。

β) 对于A-模的每个滤过归纳系统 \((M_{\alpha})_{\alpha \in I}\) ，典范同态

\[
\varinjlim \mathrm{H} ^ {i} (\operatorname{Homgr} _ {\mathrm{A}} (\mathrm{P}, \mathrm{M} _ {\alpha})) \rightarrow \mathrm{H} ^ {i} (\operatorname{Homgr} _ {\mathrm{A}} (\mathrm{P}, \varinjlim \mathrm{M} _ {\alpha}))
\]

对 \(i < n\) 是双射的，对 \(i = n\) 是单射的。

γ) 对任意右 A-模族 \((\mathbf{N}_{\alpha})_{\alpha\in I}\) ，典范同态

\[
\mathrm{H} _ {i} ((\prod_ {\alpha} \mathrm{N} _ {\alpha}) \otimes_ {\mathrm{A}} \mathrm{P}) \rightarrow \prod_ {\alpha} \mathrm{H} _ {i} (\mathrm{N} _ {\alpha} \otimes \mathrm{P})
\]

对 \(i < n\) 是双射的，对 \(i = n\) 是满射的。

（为了证明 \(\alpha\) ）蕴含 \(\beta\) ）和 \(\gamma\) ），将习题 4 应用于 \(f\) 的锥。为证明 \(\beta\) ）

或 γ) 蕴含 α)，对 n 进行归纳，应用习题 4、第 170 页习题 18 及第 174 页习题 10。)

\begin{enumerate}
\def\labelenumi{\arabic{enumi})}
\setcounter{enumi}{5}
\tightlist
\item
  设 A、B 为两个 k-代数（结合的且含单位元），M 为左 A-模，P 为左 B-模，N 为 (B, A)-双模。
\end{enumerate}

\begin{enumerate}
\def\labelenumi{\alph{enumi})}
\tightlist
\item
  证明存在（X，第 175–177 页，习题 14 至 17）两个 \(k\) -模的谱序列 \({}^{1}\mathrm{E}\) 和 \({}^{2}\mathrm{E}\) ，收敛到相同的 \(k\) -分次模，使得
\end{enumerate}

\[
{ } ^ { 1 } \mathrm{E} _ { 2 } ^ { p q } = \operatorname{Ext} _ { \mathbf { A } } ^ { p } ( \mathbf { M } , \operatorname{Ext} _ { \mathbf { B } } ^ { q } ( \mathbf { N } , \mathbf { P } ) ) ; \quad { } ^ { 2 } \mathrm{E} _ { 2 } ^ { p q } = \operatorname{Ext} _ { \mathbf { B } } ^ { p } ( \operatorname{Tor} _ { q } ^ { \mathbf { A } } ( \mathbf { N } , \mathbf { M } ) , \mathbf { P } ) .
\]

\begin{enumerate}
\def\labelenumi{\alph{enumi})}
\setcounter{enumi}{1}
\tightlist
\item
  通过复合同态 \(\gamma\) 和 \(\delta\) （习题 15，第 176 页），定义一个 \(k\) -同态
\end{enumerate}

\[
\mu : \operatorname{Ext} _ {\mathbf {A}} (\mathbf {M}, \operatorname{Hom} _ {\mathbf {B}} (\mathbf {N}, \mathbf {P})) \rightarrow \operatorname{Hom} _ {\mathbf {B}} \left(\operatorname{Tor} ^ {\mathbf {A}} (\mathbf {N}, \mathbf {M}), \mathbf {P}\right).
\]

若P是内射的，证明μ是同构。

\begin{enumerate}
\def\labelenumi{\alph{enumi})}
\setcounter{enumi}{2}
\item
  若N是平坦A-模，证明该序列 \({}^{1}E\) 收敛到Ext \(_{B}(N \otimes_{A} M, P)\) ；若 N 是投射 B-模，则序列 \({}^{2}E\) 收敛到Ext \(_{A}(M, Hom_{B}(N, P))\) .
\item
  设 \(\rho: A \to B\) 一个 \(k\) -代数。证明存在一个谱序列 E，使得 \(\mathbf{E}_2^{pq} = \operatorname{Ext}_{\mathbf{B}}^p (\operatorname{Tor}_q^A(B, M), P)\) ，收敛到 \(\operatorname{Ext}_{\mathbf{A}}(M, P)\) 以及一个谱序列 'E，使得
\end{enumerate}

\[
{ } ^ { \prime } \mathrm{E} _ { 2 } ^ { p q } = \operatorname{Ext} _ { \mathbf { B } } ^ { p } ( \mathbf { P } , \operatorname{Ext} _ { \mathbf { A } } ^ { q } ( \mathbf { B } , \mathbf { M } ) ) ,
\]

收敛到 \(\operatorname{Ext}_{\mathbf{A}}(\mathbf{P}, \mathbf{M})\) 。

\begin{enumerate}
\def\labelenumi{\arabic{enumi})}
\setcounter{enumi}{6}
\tightlist
\item
  设 C 是一个 A-复形，M 是一个 A-模。称 M 被 C 的扩张（相应地，C 被 M 的扩张）的模为分次 k-模
\end{enumerate}

\[
\mathcal {E} x t _ {\mathrm{A}} (\mathrm{C}, \mathrm{M}) = \mathrm{H} (\operatorname{Homgr} _ {\mathrm{A}} (\mathrm{C}, \mathrm{I} (\mathrm{M}))) \quad \left(\text { resp. } \mathcal {E} x t _ {\mathrm{A}} (\mathrm{M}, \mathrm{C}) = \mathrm{H} (\operatorname{Homgr} _ {\mathrm{A}} (\mathrm{L} (\mathrm{M}), \mathrm{C}))\right).
\]

若 \(f: \mathbf{C} \to \mathbf{C}'\) 是一个复形态射且 \(g: \mathbf{M} \to \mathbf{M}'\) 一个A-模的同态，我们设 \(\mathcal{E}xt_{\mathbf{A}}(f, g) = \mathrm{H}(\mathrm{Homgr}(f, \mathrm{l}(g)))\) , \(\mathcal{E}xt_{\mathbf{A}}(g, f) = \mathrm{H}(\mathrm{Homgr}(\mathrm{L}(g), f))\) .

\begin{enumerate}
\def\labelenumi{\alph{enumi})}
\item
  若 \(f: \mathbf{C} \to \mathbf{C}'\) 是一个拟同构，则 \(k\) -同态 \(\mathcal{E}xt_{\mathbf{A}}(f, 1)\) 与 \(\mathcal{E}xt_{\mathbf{A}}(1, f)\) 都是双射的。若C是投射的且在右侧有界（相应地，内射的且在左侧有界），则 \(k\) -分次模 \(\mathcal{E}xt_{\mathbf{A}}(\mathbf{C}, \mathbf{M})\) (相应地， \(\mathcal{E}xt_{\mathbf{A}}(\mathbf{M}, \mathbf{C})\) )同构于 \(\mathrm{H}(\mathrm{Homgr}_{\mathbf{A}}(\mathbf{C}, \mathbf{M}))\) (相应地， \(\mathrm{H}(\mathrm{Homgr}_{\mathbf{A}}(\mathbf{M}, \mathbf{C}))\)) 。
\item
  设 \(0 \to C' \to C \to C'' \to 0\) 是A-复形的一个正合序列，且 \(0 \to M' \to M \to M'' \to 0\) 是A-模的一个正合序列。证明存在正合序列
\end{enumerate}

\[
\begin{array}{l} \dots \to \mathcal {E} x t _ {\mathrm{A}} ^ {n} (\mathrm{C} ^ {\prime \prime}, \mathrm{M}) \to \mathcal {E} x t _ {\mathrm{A}} ^ {n} (\mathrm{C}, \mathrm{M}) \to \mathcal {E} x t _ {\mathrm{A}} ^ {n} (\mathrm{C} ^ {\prime}, \mathrm{M}) \to \mathcal {E} x t _ {\mathrm{A}} ^ {n + 1} (\mathrm{C} ^ {\prime \prime}, \mathrm{M}) \to \dots \\ \dots \to \mathcal {E} x t _ {\mathrm{A}} ^ {n} (\mathrm{C}, \mathrm{M} ^ {\prime}) \to \mathcal {E} x t _ {\mathrm{A}} ^ {n} (\mathrm{C}, \mathrm{M}) \to \mathcal {E} x t _ {\mathrm{A}} ^ {n} (\mathrm{C}, \mathrm{M} ^ {\prime \prime}) \to \mathcal {E} x t _ {\mathrm{A}} ^ {n + 1} (\mathrm{C}, \mathrm{M} ^ {\prime}) \to \dots \\ \dots \to \mathcal {E} x t _ {\mathrm{A}} ^ {n} (\mathrm{M}, \mathrm{C} ^ {\prime}) \to \mathcal {E} x t _ {\mathrm{A}} ^ {n} (\mathrm{M}, \mathrm{C}) \to \mathcal {E} x t _ {\mathrm{A}} ^ {n} (\mathrm{M}, \mathrm{C} ^ {\prime \prime}) \to \mathcal {E} x t _ {\mathrm{A}} ^ {n + 1} (\mathrm{M}, \mathrm{C} ^ {\prime}) \to \dots \\ \dots \to \mathcal {E} x t _ {\mathrm{A}} ^ {n} (\mathrm{M} ^ {\prime \prime}, \mathrm{C}) \to \mathcal {E} x t _ {\mathrm{A}} ^ {n} (\mathrm{M}, \mathrm{C}) \to \mathcal {E} x t _ {\mathrm{A}} ^ {n} (\mathrm{M} ^ {\prime}, \mathrm{C}) \to \mathcal {E} x t _ {\mathrm{A}} ^ {n + 1} (\mathrm{M} ^ {\prime \prime}, \mathrm{C}) \to \dots \end{array}
\]

\begin{enumerate}
\def\labelenumi{\alph{enumi})}
\setcounter{enumi}{2}
\tightlist
\item
  证明存在一个谱序列 ``E（X，第 175–177 页，习题 14 至 17）收敛到 \(\mathcal{E}xt_{\mathrm{A}}(\mathrm{C},\mathrm{M})\) 使得 \(E_2^{pq} = H^q (\mathrm{Ext}_A^p (\mathrm{C},\mathrm{M}))\) 若 C 上有界，证明存在一个谱序列 'E，收敛到 \(\mathcal{E}xt_{\mathrm{A}}(\mathrm{C},\mathrm{M})\) ，使得 \(E_2^{pq} = \mathrm{Ext}_A^p (\mathrm{H}_q(\mathrm{C}),\mathrm{M})\) （考虑双复形 \(\mathrm{Homgr}_{\mathrm{A}}(\mathrm{C},\mathrm{I}(\mathrm{M}))\) ）。
\end{enumerate}

\begin{enumerate}
\def\labelenumi{\arabic{enumi})}
\setcounter{enumi}{7}
\tightlist
\item
  设 P, Q 为两个 A-复形。假设 P 下有界且 Q 上有界，或者存在整数 d 使得每个 A-模都有长度为 \(\leqslant d\) 的投射分解。
\end{enumerate}

\begin{enumerate}
\def\labelenumi{\alph{enumi})}
\tightlist
\item
  证明存在两个谱序列 'E 和 ``E，收敛到同一个分次 k-模，使得
\end{enumerate}

\[
^ {\prime} \mathrm{E} _ {2} ^ {p q} = \mathrm{H} ^ {p} (\mathrm{Ext} _ {\mathrm{A}} ^ {q} (\mathrm{P}, \mathrm{Q})) \quad^ {\prime \prime} \mathrm{E} _ {2} ^ {p q} = \bigoplus_ {q ^ {\prime} + q ^ {\prime \prime} = q} \mathrm{Ext} _ {\mathrm{A}} ^ {p} (\mathrm{H} _ {q ^ {\prime}} (\mathrm{P}), \mathrm{H} ^ {q ^ {\prime \prime}} (\mathrm{Q}))
\]

\((\text{或} \operatorname{Ext}_{\mathbf{A}}^{q}(\mathbf{P}, \mathbf{Q}) \text{ 表示与双复形相关联的复形 } \oplus \operatorname{Ext}_{\mathbf{A}}^{q}(\mathbf{P}_{r}, \mathbf{Q}^{s}))\) .

\begin{enumerate}
\def\labelenumi{\alph{enumi})}
\setcounter{enumi}{1}
\tightlist
\item
  若 P 是投射的或 Q 是内射的，则谱序列 ``E 收敛到 H(Homgr \(_A\) (P, Q))。若 H(P) 是投射的或 H(Q) 是内射的，则谱序列 'E 收敛到 Homgr \(_A\) (H(P), H(Q))。
\end{enumerate}

